% Generated draft: requires mathematical and source review.
\documentclass[11pt]{article}
\usepackage[margin=25mm]{geometry}
\usepackage{fontspec}
\setmainfont{Latin Modern Roman}
\usepackage{amsmath,amssymb,mathtools,unicode-math}
\setmathfont{Latin Modern Math}
\usepackage{enumitem,needspace,xurl,hyperref}
\setlength{\parindent}{0pt}
\setlength{\parskip}{5pt}
\setlength{\emergencystretch}{3em}
\newcommand{\N}{\mathbb{N}}
\newcommand{\Z}{\mathbb{Z}}
\newcommand{\Q}{\mathbb{Q}}
\newcommand{\R}{\mathbb{R}}
\newcommand{\C}{\mathbb{C}}
\newcommand{\F}{\mathbb{F}}
\newcommand{\A}{\mathbb{A}}
\newcommand{\PP}{\mathbb{P}}
\newcommand{\E}{\mathbb{E}}
\newcommand{\eps}{\varepsilon}
\newcommand{\dd}{\,\mathrm{d}}
\begin{document}
\section*{UnsolvedMath ID: 30001051}
\textbf{Source problem code:} catalogue entry
\subsection{Definitions and mathematical statement}
A hyperbolic knot manifold is understood here as a compact connected orientable three-manifold $M$ with torus boundary whose interior has a complete finite-volume hyperbolic metric. It is not restricted to knot complements in $S^3$. A slope is an unoriented isotopy class of essential simple closed curves on $\partial M$. The Dehn filling $M(\alpha)$ attaches a solid torus with meridian identified to slope $\alpha$.

Let $E(M)$ consist of the slopes for which $M(\alpha)$ is not hyperbolic. For slopes $\alpha,\beta$, their distance $\Delta(\alpha,\beta)$ is their minimal geometric intersection number, equal to $|ps-qr|$ for primitive coordinate pairs $(p,q),(r,s)$ in a common homology basis. Put
\[
\Delta(E(M))=\max_{\alpha,\beta\in E(M)}\Delta(\alpha,\beta),
\]
with value zero for fewer than two slopes.

Retain both levels of the source's proposed bounds:
\[
\forall M,\qquad |E(M)|\leq10,\qquad\Delta(E(M))\leq8;
\]
and
\[
M\notin\{M_1,M_2,M_3,M_4\}\quad\Longrightarrow\quad
|E(M)|\leq7,\qquad\Delta(E(M))\leq5.
\]
The supplied assessment reports that Lackenby--Meyerhoff proved the first pair; they are preserved as historically posed statements with that solved-part qualification. The sharper exclusion-dependent pair remains the additional target in the supplied record.

The four named exceptional manifold types $M_1,\ldots,M_4$ are not identified by constructions, surgery descriptions or census names. Their intended types must be recovered before the stronger statement has a fully specified exclusion. No particular four manifolds or restriction to $S^3$ knot complements is guessed.
\subsection{Short English statement}
Retain the universal exceptional-filling bounds $10$ and $8$, reported proved by Lackenby--Meyerhoff, and the sharper proposed bounds $7$ and $5$ outside four named manifold types. The four exceptions are not defined in the extract.
\subsection{Sources}
\begin{itemize}
\item Supplied report: \emph{Topologie} (2009), Oberwolfach Reports, DOI 10.4171/owr/2008/43, pp. 2419--2476; proposed year 2008. \url{https://doi.org/10.4171/owr/2008/43}.
\item Supplied Lackenby--Meyerhoff reference: \url{https://arxiv.org/abs/0808.1176}.
\item Supplied Gordon retrospective: \url{https://celebratio.org/Gordon_C/article/1165/}.
\item Local UnsolvedMath registry, permanent ID 30001051, catalogue entry; created 25 August 2026, updated 16 September 2026. Its corrected\_verified note reports comparison with the official report. Literature checked at \texttt{2026-08-21}: the assessment is supplied dated curation, not independently verified. It labels the record partially solved, reports proof of $|E(M)|\leq10$ and diameter at most $8$, and says the sharper bounds outside the four exceptions remain partly unresolved.
\end{itemize}
\end{document}
